\chapter{The Ledger's Permissions}
\label{chap:ledger-permissions}

% EPISODE COVERAGE: Episode 2 -- ADMISSIBLE through NUMERIC, Sample,
% ObservationProcess, BINARY, REPEATABLE, ComputationalProcess.closure.
% KEY CONCEPT: permissions as the rules that turn number forms into a ledger.

A threshold is not only a gate. A gate suggests a door, a hinge, a yes or no.
But the instruments in this chapter need something more precise than a door.
They need a way to decide what may become part of a record, how fast that
record may grow, whether the next entry actually moves the record forward, and
when further comparison has gone quiet. A threshold is the place where a
physical process becomes a ledger process.

This is why thresholds have to be set before the record begins. A threshold
chosen after the mark arrives is not a threshold. It is an excuse. The line on
the glass, the voltage across the tube, the exposure range of the paper, the
trial protocol signed before enrollment: each one is a promise that the
instrument will not change its rule after seeing what came through. A record
can tolerate noise. It can tolerate missed events. It cannot tolerate a moving
standard and still call itself a measurement.

Consider a Geiger counter. A particle may pass near the tube and still leave
no entry. The counter records only an ionization strong enough to trigger the
avalanche. Below that avalanche, something may have happened in the tube, but
the instrument has no mark it is allowed to keep. The record begins only when
the event clears the entry threshold.

The counter has another threshold hidden inside its pace. After a discharge,
the tube needs a short recovery interval. During that dead time, a second
particle may arrive with more than enough energy to matter physically. The
instrument still cannot count it. The record is bounded by the counter's
capacity to recover. It may grow one way, but not arbitrarily fast. The
instrument does not merely ask whether an event occurred. It asks whether the
event occurred in a window where it could be distinguished from the event just
recorded.

The record also has a direction. Counts are timestamped. A later count can
extend the run; a count whose timestamp falls behind the previous count cannot
be inserted as if nothing happened. The counter's ledger accumulates one way.
It may miss events. It may saturate. It may have uncertainty. But it cannot
honestly move backward through its own record.

Finally, a radiation measurement does not close because the universe has
stopped emitting particles. It closes because the comparison has settled
enough for the instrument's purpose. A count over one dwell time is compared
with a background count over the same dwell time. More counting may reduce the
uncertainty, but at some point the ratio of signal to background stops changing
in a useful way. The physicist sets down the detector when the remaining
difference is smaller than the decision the record was built to support.
Before that moment, the tube must also be calibrated against its own ordinary
silence. Background is not nothing. It is the instrument's watched absence, the
rate at which the world enters the counter when no source is being tested.
Only against that watched absence can the source count become a reading rather
than a pile of clicks.

A photographic emulsion has the same four thresholds, but now they appear in
light and chemistry. A silver halide grain does not become part of the image
because one imagines light nearby. It needs enough photons to initiate the
latent change that development can reveal. Below that threshold, light may
have brushed the material, but the grain will not enter the record.

The emulsion also bounds the growth of the record. Its tonal range is finite.
At the dark end, small differences disappear into the base. At the bright end,
the material saturates. More exposure does not produce a more precise reading;
it collapses many differences into the same dense patch. The image can grow in
detail only inside the range the material can distinguish.

Development gives the record its direction. Expose, develop, stop, fix: the
order is not a preference. If the fixer arrives before the developer, the
latent image is destroyed. If the stop bath is skipped, development keeps
moving past the useful point. The print is not a pile of chemical states. It
is a progression whose steps must occur in the order the material permits.

The print is lifted from the tray when further development no longer changes
what the eye can use. It may still be possible to leave the paper longer. That
is not the point. The comparison between the emerging image and the intended
record has settled. The next moment in the bath no longer earns a new
distinction.
The darkroom teaches another part of the same discipline: the threshold must
belong to the material, not to the wish of the photographer. Paper has a grade.
Developer has a strength. The safelight is safe only because its wavelength
falls beneath what the emulsion will register within the working time. The
record is possible because the room itself has been arranged around what the
paper can and cannot admit.

The Geiger tube and the photographic paper are very different instruments. Yet
each has the same hidden grammar. Something must clear an entry threshold. The
record must be bounded in its growth. Its steps must move forward rather than
rewriting themselves. It must know when further comparison has fallen silent.
These are not four separate phenomena. They are four faces of one problem:
what may become part of the record.

The number forms are built. Each can be compared, ordered, and accumulated.
What the ledger cannot yet do is decide which readings may enter, how fast
they may grow, which progressions are legitimate, and when two settling paths
represent the same measurement. These are not properties of the numbers. They
are permissions the ledger must be given. This chapter installs them.

Each permission is named for what it authorizes, and each is tied to a number
form the previous chapter built.

\begin{center}
\begin{tabular}{@{}p{0.24\textwidth}p{0.22\textwidth}p{0.42\textwidth}@{}}
\lean{ADMISSIBLE} & \lean{Natural} &
what may enter as evidence \\
\lean{COUNTABLE} & \lean{Rational} &
how pieces are bounded and numbered \\
\lean{ENCODED} & \lean{Sequence} &
how readings become forward progressions \\
\lean{RESIDUE} & \lean{Limit} &
when settling paths no longer differ usefully
\end{tabular}
\end{center}

In the courtroom metaphor: admission, numbering, comparison, verdict.

\section{The Chain of Admissibility}
\label{se:chain-of-admissibility}

The number ladder is now visible: \lean{Natural}, \lean{Rational},
\lean{Sequence}, and \lean{Limit}. This chapter does not introduce new number
types. It authorizes the use of those types. The previous chapter built the
forms an instrument can hold. This chapter installs the promises that let the
ledger use them. Each promise names a permission the ledger did not previously
have.

\lean{CountingProcess} packages a carrier, a \lean{Natural}, and an iterator:
the fixture required before the ledger can treat receipt-stacks as a count.

\begin{leancode}
class ADMISSIBLE\\
\quad (Value: Type)\\
\quad (Carrier: CarrierProcess Value)\\
\quad [DISTINGUISHABLE Value Carrier]\\
\quad where\\
\quad counting\_process : CountingProcess Value Carrier\\
\quad admissible? : Number -> Number -> Prop :=\\
\quad\quad fun stimulus threshold =>\\
\quad\quad match stimulus, threshold with\\
\quad\quad | .zero \_, \_ => True\\
\quad\quad | .one \_ \_, .zero \_ => False\\
\quad\quad | .one \_ n1', .one p2 n2' =>\\
\quad\quad\quad p2.truth \(\vee\) n1' <= n2'
\end{leancode}

The stimulus and threshold are \lean{Number}s, the receipt-stacks from Chapter
1. A zero stimulus is always admitted. A nonzero stimulus against a zero
threshold is rejected. When both sides are nonzero, the threshold's fact may
already let the stimulus through, or the stimulus tail must fit below the
threshold tail. This is admission: a distinguished mark becomes evidence only
if it clears the gate.

The move to \lean{Natural} is the move from a single receipt to a directed
count. \lean{ADMISSIBLE} is the promise that this move is legitimate for the
carrier at hand. It does not say that every possible mark belongs in the
ledger. It says that, once a mark has been distinguished, there is a gate by
which it may be tested against the current threshold. The natural number
appears here as the running stack made possible by that gate. Without
admission, a pile of receipts is only a pile. With admission, the pile becomes
a count whose next step has been allowed to enter.

This is the class pattern for the permission chain. Each later class carries
the same \lean{Value} and \lean{Carrier} parameters, inherits the promises
already installed, and adds one new constraint to the stack. From here on, the
code blocks show only the new constraint and the predicate that does the
conceptual work. The inherited parameter block and process field are omitted.

\lean{IndexingProcess} packages a \lean{Rational} origin with the counting
process that advances it.

\begin{leancode}
[ADMISSIBLE Value Carrier]\\
bounded? : Rational -> Rational -> Prop :=\\
\quad\quad fun stimulus threshold =>\\
\quad\quad match stimulus, threshold with\\
\quad\quad | .zero \_, \_ => True\\
\quad\quad | .number \_ \_ \_, .zero \_ => False\\
\quad\quad | .number p1 i1 \_, .number p2 i2 \_ =>\\
\quad\quad\quad (p1.truth = p2.truth \(\wedge\) i1 <= i2) \(\vee\)\\
\quad\quad\quad (p1.truth \(\neq\) p2.truth \(\wedge\) i2 <= i1)
\end{leancode}

The move to \lean{Rational} is the move from whole ticks to comparable pieces.
\lean{bounded?} asks whether one indexed piece falls within another. When the
truths agree, the indexes are read forward. When the truths differ, the order
reverses. Numbering is also a bound: the record may grow only as fast as the
instrument can distinguish what is being added.

\lean{COUNTABLE} therefore authorizes more than incrementation. It says that a
new piece can be placed against an origin and compared with a threshold of the
same kind. The familiar rational intuition is useful here: a ratio is not
merely a bigger count, but a count read against another count. The class keeps
that intuition inside the ledger's discipline. It does not import division.
It asks whether the proposed indexed reading remains within a bounded
comparison. The contravariant branch is essential, because an instrument
sometimes learns a quantity by reading the same relation from the opposite
side. A receding Doppler return and an approaching one are both readings, but
their order is not interpreted by the same branch. The promise of
\lean{COUNTABLE} is that the ledger can make that branch choice without
forgetting the receipt structure that made counting possible in the first
place.

\begin{gphenom}{Berkeley-Galileo}
\label{ph:berkeley-galileo}
\GPhObs Berkeley argued in 1734 that Galileo's law of free fall implies an
instantaneous infinite velocity at the moment of release, which Berkeley
considered an absurdity that undermined the calculus.
\GPhCon The rate at which new facts can be added to the ledger is bounded by
the instrument's resolution. The derivative of the ledger is bounded
by the admissibility rate.
\GPhSeq \lean{COUNTABLE}'s \lean{bounded?} predicate is this constraint made
formal: no sequence of events may exceed the instrument's capacity to
distinguish them. The ledger grows at most as fast as the instrument can witness.
\end{gphenom}

\lean{LimitProcess} packages a \lean{Sequence}, a rational limit value, and an
iterator for advancing the ordered progression.

\begin{leancode}
[COUNTABLE Value Carrier]\\
encoding? : Sequence -> Sequence -> Prop :=\\
\quad\quad fun s1 s2 =>\\
\quad\quad match s1, s2 with\\
\quad\quad | .nil f1, .nil f2 => f1 = f2\\
\quad\quad | .nil \_, .index \_ \_ \_ => True\\
\quad\quad | .index \_ \_ \_, .nil \_ => False\\
\quad\quad | .index f1 r1 i1, .index f2 r2 i2 =>\\
\quad\quad\quad (f1 = f2) \(\wedge\) r1 <= r2 \(\wedge\) i1 < i2
\end{leancode}

A \lean{Sequence} is what happens when indexed pieces appear in an order that
can move forward. The strict inequality \lean{i1 < i2} is the formal statement
that the comparison must have advanced. \lean{ENCODED} is where the record
becomes more than a list of bounded values. The chain acquires an arrow.

\lean{ENCODED} is the promise attached to \lean{Sequence}. A sequence is not
only a collection of rational readings. It is a record in which one indexed
piece follows another under a comparison the ledger can check. The class
requires a \lean{LimitProcess}, but the immediate work of the predicate is
directional. Nil may begin the progression. A nonempty entry cannot be below
nil. Two indexed entries agree only when their facts agree, their rational
readings are ordered, and the later index is strictly later. This is how a
time-ordered instrument record differs from a bag of samples. The promise does
not yet say that the record has settled. It says that the record has been
encoded as progress, so later claims may ask whether the progress converges.

\lean{CauchyProcess} packages a \lean{Limit} with a current rational value and
an accumulation.

\begin{leancode}
[ENCODED Value Carrier]\\
representative? : Limit -> Limit -> Prop :=\\
\quad\quad fun a b =>\\
\quad\quad match a, b with\\
\quad\quad | .nil f1, .nil f2 => f1 = f2\\
\quad\quad | .nil \_, .index \_ \_ \_ => True\\
\quad\quad | .index \_ \_ \_, .nil \_ => False\\
\quad\quad | .index f1 s1 l1, .index f2 s2 l2 =>\\
\quad\quad\quad (f1 = f2) \(\wedge\) s1 <= s2 \(\wedge\) l1 < l2
\end{leancode}

A \lean{Limit} is not a completed real number here. It is the shape of a
settling record. \lean{representative?} declares when two limit paths can no
longer be told apart by admissible comparison. This is equivalence by
comparison, not equality by identity. The ledger does not require two
measurement paths to be the same path. It requires that no relevant difference
remain. That is the courtroom verdict.

\lean{RESIDUE} is the promise attached to \lean{Limit}. It does not create a
new continuum for the ledger to inhabit. It gives the ledger a way to say that
two accumulated progressions have the same relevant remainder. The familiar
mathematical word is convergence, but the device phrases the point as a
representative test. A nil limit may be compared with another nil limit by
fact equality. A nil limit sits below a nonempty path. Two nonempty paths
compare by fact, by sequence order, and by a further limit index. The structure
is recursive because settling is not a single final glance. It is the
continued failure of further admissible comparison to expose a difference that
matters to the record.

The wheel-speed sensor has now gained more than a pulse count. It can recognize
a tooth passing, admit that tooth into the record, number it as a bounded
piece, compare recordings as ordered progressions, and declare when two
settling paths represent the same measurement. The number thread that began
with a single witnessed mark now runs through four types, four processes, and
four class promises. The ledger is ready to hold a sample.

\section{The First Multi-Fact Structure}
\label{se:first-multi-fact}

The preceding sections built enough numerical structure to say what counts,
what is bounded, what progresses, and what has settled. But a measurement is
not only a settled path. It is an event in which a path is taken as a sample.
That requires a structure with room for more than one fact.

\begin{leancode}
inductive Sample where\\
\quad | initial\_condition : Fact -> Limit -> Sample\\
\quad | signal\_response :\\
\quad\quad Fact -> Limit -> Fact -> Limit -> Sample -> Sample
\end{leancode}

\lean{Sample} is the first multi-fact structure in the tower. The initial case
stores the fact and limit from which the instrument begins. The response case
stores a first fact and limit, a second fact and limit, and the sample that
came before. The record is therefore not a bag of readings. It is an ordered
accumulation of before-and-after structure.

The order on samples repeats the pattern one more time.

\begin{leancode}
namespace Sample\\
def le : Sample -> Sample -> Prop\\
\quad | .initial\_condition f1 l1, .initial\_condition f2 l2 =>\\
\quad\quad (f1.truth = f2.truth \(\wedge\) l1 <= l2) \(\vee\)\\
\quad\quad (f1.truth \(\neq\) f2.truth \(\wedge\) l2 <= l1)\\
\quad | .signal\_response f1 s1 \_ r1 \_,\\
\quad\quad .signal\_response f2 s2 \_ r2 \_ =>\\
\quad\quad (f1.truth = f2.truth \(\wedge\) r1 <= r2) \(\vee\)\\
\quad\quad (f1.truth \(\neq\) f2.truth \(\wedge\) s1 <= s2)\\
\quad | .initial\_condition f1 l1,\\
\quad\quad .signal\_response f2 l2 \_ \_ \_ =>\\
\quad\quad (f1.truth = f2.truth \(\wedge\) l1 <= l2) \(\vee\)\\
\quad\quad (f1.truth \(\neq\) f2.truth \(\wedge\) l2 <= l1)\\
\quad | .signal\_response \_ \_ \_ \_ \_, .initial\_condition \_ \_ =>\\
\quad\quad False
\end{leancode}

The important fact is not the detail of every case. It is the shape. Initial
conditions compare by their limits. Signal responses compare by the relation
between what came in and what came back. A later response cannot be moved back
under an initial condition. Once a response has been recorded, the ledger has
crossed a boundary.

\lean{Sample} and \lean{Limit} therefore split observation into two roles. The
\lean{Limit} is the thing observed, the settling path. The \lean{Sample} is the
act of taking that path into the record. This is where the ledger first looks
like an experiment rather than only a number system.

\begin{gphenom}{da Vinci-Coulomb}
\label{ph:da-vinci-coulomb}
\GPhObs Da Vinci observed that a block on a rough surface resists motion
until the applied force exceeds a threshold. Coulomb formalized this in 1781:
static friction holds until the threshold is crossed, then yields.
\GPhCon A proposed continuation is not admissible until it exceeds the ledger's
static friction threshold -- the accumulated inconsistency required to force
a new entry. Below the threshold, the ledger absorbs the proposal without yielding.
\GPhSeq The \lean{Sample} structure models this exactly: a new response can be
appended to the record only when the incoming mark has cleared the
admissibility threshold. Below that threshold, the record does not grow.
Measurement is not continuous. It is a series of threshold crossings.
\end{gphenom}

\section{The Clock and Its Complement}
\label{se:clock-and-complement}

Episode 2 begins by asking what it means to observe the sample just built. The
answer is not a new number. It is a process with a before and an after.

The parameter block and constraint stack are the same as in the preceding
classes and will not be shown again for structures. The fields are the new
fixture.

\begin{leancode}
\quad cauchy\_process : CauchyProcess Value Carrier\\
\quad before : Limit\\
\quad after : Limit\\
\quad relative\_variance : Prop
\end{leancode}

\lean{ObservationProcess} carries the whole tower forward. It requires
distinction, admission, counting, encoding, and residue before it is allowed to
name the two sides of an observation. The fields \lean{before} and
\lean{after} are both \lean{Limit}s because observation happens at the level
where readings have settled enough to be compared.

This is the structure pattern for the rest of the chapter. Later structure
blocks show only the fields that change.

The new predicate is \lean{relative\_variance}. In the Lean file it asks
whether the facts attached to the before and after limits agree. If they do,
the observation is covariant. If they do not, the observation is
contravariant. The old truth-parity rule has become a clock.

The class \lean{BINARY} packages that clock.
\lean{BINARY} breaks the preceding class pattern. Where the permission classes
carried one process and one predicate, this class carries the observed unit
itself: zero, one, the bit between them, and a distinction test.

\begin{leancode}
\quad zero : Limit\\
\quad one : Limit\\
\quad bit : Sample\\
\quad different? : Sample -> Sample -> Prop
\end{leancode}

\lean{BINARY} does not introduce a binary numeral system in the usual sense.
It introduces the minimum unit of observed difference: a zero limit, a one
limit, and a sample that can tick between them. A clock says true, false, true,
false. Its complement says false, true, false, true. Together they tile the
same distinction from opposite phases.

\begin{sidebar}{Why binary?}
\label{sb:why-binary}
A single distinction is the minimum unit of information. The ledger does not
need an alphabet yet. It needs one stable difference, and then the complement
of that difference. From that pair, it can begin to count observations.
\end{sidebar}

This is why the tower began with \lean{DISTINGUISHABLE}. A bit is not a tiny
number floating in isolation. It is a witnessed distinction that survived the
whole tower: admitted, bounded, encoded, reduced to residue, and observed as a
before-after pair.

The same pattern then moves one level higher, from samples to trials.

\lean{Trial} is to \lean{Sample} what \lean{Sample} is to \lean{Limit}: the
same two-constructor pattern, now carrying samples instead of limits. A trial
turns an observed bit into something that can be repeated.

\section{Closure}
\label{se:closure-first-appearance}

Once the ledger can observe a bit, it can repeat the observation. Episode 2
does this by building repeatable processes and studies. The repeatable process
packages an observation process, a sample stimulus, and an expected trial.

\begin{leancode}
\quad stimulus : Sample\\
\quad expectation : Trial\\
\quad iterate : Trial -> Trial
\end{leancode}

A \lean{RepeatableProcess} keeps the observation machinery together with the
stimulus and expected trial that make one run repeatable.

\begin{leancode}
[BINARY Value Carrier]\\
typical\_response : Trial -> Trial -> Prop
\end{leancode}

\lean{REPEATABLE} requires that process and a typical-response predicate: the
formal version of an experimental protocol, saying that given a stimulus trial,
this is the trial the process predicts.

The ledger also needs an open account.

\begin{leancode}
inductive Study where\\
\quad | hypothesis : Fact -> Study\\
\quad | data : Fact -> Trial -> Study -> Study
\end{leancode}

A \lean{Study} is the ledger's open account: either a hypothesis or a
witnessed fact attached to a trial and the study that preceded it. Only now can
the chapter name the first explicit computational closure.

\begin{leancode}
\quad output : Option Study\\
\quad close : Option Study -> Study\\
\quad closure : Study -> Study
\end{leancode}

\lean{ComputationalProcess} is the first place where the ledger says what it
does with an open record. The field \lean{output} may contain a completed
\lean{Study}, or it may contain nothing. The function \lean{close} turns either
case into a study. If there is an output, it returns it. If there is none, it
returns the error-code study supplied by the process.

Then \lean{closure} turns a study into another study. A hypothesis becomes data
by attaching the expected trial and the closed output. Existing data becomes
new data by iterating the trial and closing the output again. In compiler
language, this is what it feels like when all holes are filled, all local
obligations are resolved, and the elaborator has run to completion. In
measurement language, the proposed continuation has become consistent with the
record.

\lean{NUMERIC} wraps that closure in the next class promise.

\begin{leancode}
[REPEATABLE Value Carrier]\\
carrier : Study\\
lambda : Option Study -> Study\\
related : Study -> Study -> Prop
\end{leancode}

\lean{NUMERIC} does not say that the carrier is a number in the ordinary
sense. It says that the carrier can now be used numerically. There is a study
that can carry it, a lambda that closes optional output, and a relation that
orders one study against another. The tower has not imported arithmetic. It has
built enough ledger discipline for arithmetic behavior to have a place to sit.

This is the chapter's first full closure. The compiler has agreed to a sequence
of promises: distinguish, admit, count, bound, encode, reduce, sample, observe,
repeat, and close. The ledger is not yet the whole instrument. It is not yet
geometry, calculus, or physics. But it now has a closed account. A mark can
enter, become evidence, become a count, become a bounded piece, become a
progression, become a residue, become a sample, become a bit, and return as a
closed study.

This is the substrate the next chapter needs. Once a closed study can carry
counts, ratios, progressions, and settling values, a signal can arrive as more
than noise. It can be tested against the boundary of what the ledger can decide.
The next chapter asks what happens when a closed ledger meets undecidable
programs, noisy thresholds, and the measurement boundary.

\section*{Coda}
\markboth{CODA}{CODA}

A child's string game begins with almost nothing: one loop, two hands, and
tension. The loop does not contain a ladder, a cradle, a diamond, and a net as
separate things. It becomes those figures by being held differently. One finger
lifts a strand. Another passes under. A crossing that was slack becomes
load-bearing. The string has not changed substance. Its obligations have
changed.

Chapter 3 has been playing that game with permissions. The ladder was already
there: count, piece, path, settling value. But a loose ladder is not yet an
instrument. A strand must be admitted before it can enter the figure. A crossing
must be numbered before the next hand can find it. A passage must move forward
rather than collapse into the shape already held. A residue must settle before
the figure can be lifted.

That is why the words in this chapter are so plain. Admission, counting,
encoding, residue: they are the holds by which the ladder becomes a usable
figure. Then the figure acquires hands. A sample holds a settled path as an
event. An observation gives that event a before and an after. A bit gives the
ledger one stable difference and its complement. A trial asks whether the same
difference can be produced again. A study keeps the account open until closure
turns it into something the compiler can carry forward.

The number ladder and the permissions are not separate inventions. They are one
movement in two acts. The ladder builds the forms. The permissions say what the
forms may become. Without the ladder, the permissions have nothing to authorize.
Without the permissions, the ladder is a beautiful structure no instrument can
use.

Closure is the moment the figure can be lifted from the hands without losing
its form. The compiler has not believed the metaphor. It has checked the
crossings. Once the figure can be lifted, the next question is what motions of
the figure are allowed.
